\section{Limitations and future work}

\subsection{Limitations of the measurements}

\begin{itemize}
  \item \textbf{It is precision, not accuracy.} There was no reference of known coordinates and no RTK
        receiver. A constant error that affects all fixes equally is not seen in the scatter.
  \item \textbf{It was not open sky.} The cold starts were done in a neighborhood street between houses and
        the static test from a window with a tree and cables nearby. The height of the obstacles, the
        percentage of visible sky and the radio interference were not recorded.
  \item \textbf{The static test was short.} About 9.5 minutes of useful data remained, not the 15 planned,
        split into two runs. The two runs gave different results and their means are \dStaticRunSeparation~m
        apart, so a longer test at another time could give different numbers.
  \item \textbf{The walks are transcribed, not captured.} The phone app did not save the file, so the figures come
        from a transcription of the copied text. The walk in the open lasted \dWalkSeconds~s and the accuracy it gives is
        the receiver's estimate, not the error against a known point. Of the 15-minute static test of 20 November only the
        last minute was kept, so there is no CEP from it.
  \item \textbf{One board, one place, one day.} None of this was repeated with another board, at another place
        or at another time, which is when the visible constellation changes.
  \item \textbf{The three cold-start times were timed by hand}, with no console log, with an uncertainty of
        about 2~s, and without noting the exact weather of each one.
  \item \textbf{The satellite list does not measure the signal.} It counts satellites that appear in the
        receiver's response, not satellites with a usable signal nor used in the fix, because the firmware does
        not log signal strength.
  \item \textbf{The fence time was measured with the kernel counter}, not with an oscilloscope on a pin, and
        with regular polygons of 500~m radius, not with a real fence.
  \item \textbf{A platform different from the product's.} Everything was measured on the \dSoc{}, not on the
        ESP32-S3 nor with the u-blox receiver of the PCB, which are the targets of SR-11 and SR-12.
  \item \textbf{An unverified inference:} that the cost per vertex of the fence comes from running the
        double-precision operations in software.
\end{itemize}

\subsection{Design decisions that could change}

The 60~m square fence, the 20~m distance for pure green, the 50~m accuracy limit to accept a fix and the 10~s
to declare the fix lost are decisions of the PoC, not requirements of the Technical Report. The fence is
generated around the first fix and is not saved: a reboot rebuilds it at another place.

\subsection{Future work}

\begin{enumerate}
  \item \textbf{An open-sky accuracy test}, of at least 15 to 30 minutes, repeated at another time, and with a
        reference of known coordinates or an RTK receiver, so that accuracy can be discussed.
  \item \textbf{Signal diagnostics:} log the strength of each satellite, how many are used in the fix and which
        satellites have stored data and until when, to explain the loss of fix and the slow start, and to decide whether fixes whose accuracy worsens quickly
        should be discarded (today up to 50~m are accepted).
  \item \textbf{A user-defined polygon} (the fence of a farm), instead of a square around the first fix. It is
        planned as revision 1.1 of this report, storing the vertices in non-volatile memory and with a way to
        create them that can be carried to the final PCB.
  \item \textbf{Speed up the fence evaluation}, for example with single-precision or fixed-point arithmetic,
        which would allow polygons with more vertices within the 50~ms.
  \item \textbf{Port the firmware to the PCB:} a second GNSS adapter for the u-blox receiver, the board
        definition with the ESP32-S3, and repeat the measurements there, including the fence time with an
        oscilloscope.
  \item \textbf{A second report} with this same template for the hardware of the collar, with its own
        measurements.
\end{enumerate}
